\documentclass[12pt]{article}
\usepackage[utf8]{inputenc}
\usepackage[T1]{fontenc}
\usepackage{hyperref}
\title{Uspanteko Language Resources for K-12 Curriculum Development}
\author{Compiled by Claude Sonnet 5 (claude-sonnet-5) for Andr\'as Kornai}
\date{}
\begin{document}
\maketitle

\section{Introduction}
Uspanteko (ISO 639-3: usp) is a Mayan language of Guatemala (Quich\'e department). It is one of the smallest Mayan languages by number of speakers and is considered at risk. The Academia de Lenguas Mayas de Guatemala (ALMG) has an Uspanteko linguistic community.

\section{Primary Educational Resources}
No textbook or primer was found directly in the search results. The ALMG Uspanteko linguistic community released an educational video in May 2025 promoting learning of Uspanteko (occupations and professions vocabulary):
\begin{itemize}
\item \url{https://agn.gt/conozca-las-ocupaciones-y-profesiones-en-idioma-uspanteko/}
\item \url{https://agn.gt/?p=440461}
\end{itemize}

\section{Linguistic and Reference Materials}
\begin{itemize}
\item Wikipedia overview: \url{https://en.wikipedia.org/wiki/Uspantek_language}
\item Telma Can Pixabaj (2007), \emph{Gram\'atica Uspanteka}, OKMA, Guatemala (cited in a Glottolog record; no URL seen)
\end{itemize}
Other works cited in search summaries but without a link found: Tuyuc Sucuc, \emph{Vocabulario uspanteko} (ALMG, 2001); M\'endez and Angel, \emph{Diccionario biling\"ue uspanteko-espa\~nol} (Cholsamaj, 2007). No online access was found for these.

\section{Religious and Community Resources}
Nothing specific was found in search results.

\section{Digital and Audio Resources}
\begin{itemize}
\item Endangered Languages Archive, ``Two Mayan Languages: Uspanteko and Sakapulteko'': \url{https://eldp.access.preservica.com/dk0028/}
\item Max Planck numerals page for Uspanteko: \url{https://lingweb.eva.mpg.de/channumerals/Uspanteko.htm}
\item G\"ottingen Maya project page: \url{https://spw.uni-goettingen.de/projects/maya/lng-USP.php}
\end{itemize}

\section{Curriculum Development Context}
Existing reference works (grammar, two dictionaries) and an archived text corpus provide a base. Online, ready-to-use classroom materials were not found; coordination with the ALMG community would be necessary.
\end{document}